\chapter[Universal Security Bounds]{Universal Security Bounds}
\section{Introduction}
Chapters 14--17 developed two complementary theories on the WIS
manifold:
\begin{itemize}
\item a \textbf{smooth theory} for differentiable information functionals;
\item a \textbf{piecewise-smooth theory} for ranking-based security functionals.
\end{itemize}
This chapter unifies both perspectives into a single collection of
geometric security bounds expressed in terms of the Universal Optimality
Defect (UOD).
Throughout this chapter,
$\mathcal D(P,Q)=\|F(P)-F(Q)\|^2$
denotes the two-point Universal Optimality Defect.
\section{Smooth Security Bounds}
The results of Chapter 14 immediately imply the following general
theorem.
\begin{theorem}[Universal Smooth Bound]
\label{ch:thm:18-1}
Let $S:\mathcal P\rightarrow\mathbb R$ be a \(C^{1}\) functional satisfying
$\|\nabla_gS\|_g\le L.$
Then
\[ \boxed{
|S(P)-S(Q)|
\le
L\sqrt{\mathcal D(P,Q)}.
} \]
If \(S\in C^{2}\) and
$\nabla_gS(P^{\star})=0,$
then locally
\[ \boxed{
|S(P)-S(P^\star)|
\le
C\,\mathcal D(P,P^\star).
} \]
These estimates depend only on the WIS geometry and not on the specific
functional.
\end{theorem}
\section{Piecewise-Smooth Security Bounds}
For rank-based functionals, the previous theorem applies on each
regularity interval identified in Chapter 17.
\begin{theorem}[Local Piecewise Bound]
\label{ch:thm:18-2}
Let \(S\) satisfy the hypotheses of Theorem 17.1.
If two posterior distributions belong to the same regularity interval,
then
$\|S(P)-S(Q)\| \le L_S \sqrt{\mathcal D(P,Q)}.$
If the local representative is twice continuously differentiable and has
a stationary point,
\[ \|S(P)-S(P^{\star})\| \le C_S
\mathcal D(P,P^{\star}). \]
\end{theorem}
\section{A Unified Classification}
Every information or security functional considered in this appendix
falls into one of two categories.
\[
\begin{array}{lll}
\textbf{Functional} & \textbf{Regularity} & \textbf{Applicable Theory} \\
\hline
\text{Shannon entropy} & \text{Smooth} & \text{Chapter 14} \\
\text{R\'enyi entropy} & \text{Smooth} & \text{Chapter 14} \\
\text{Tsallis entropy} & \text{Smooth} & \text{Chapter 14} \\
\text{Collision entropy} & \text{Smooth} & \text{Chapter 14} \\
\text{Jensen--Shannon divergence} & \text{Smooth} & \text{Chapters 14--15} \\
\text{KL divergence} & \text{Smooth/Bregman} & \text{Chapters 14--15} \\
\text{Min-Entropy} & \text{Piecewise smooth} & \text{Chapter 17} \\
\text{Bayes Vulnerability} & \text{Piecewise smooth} & \text{Chapter 17} \\
\text{Guessing Entropy} & \text{Piecewise smooth} & \text{Chapter 17} \\
\text{Top-\(k\) success probability} & \text{Piecewise smooth} & \text{Chapter 17}
\end{array}
\]
Thus no separate geometric theory is required for each individual
functional.
\section{Unified Stability Principle}
The previous results can be summarized by a single principle.
\textbf{Principle.}
Every security functional considered in this appendix admits a geometric
stability estimate controlled by the Universal Optimality Defect.
Specifically,
\begin{itemize}
\item globally for smooth functionals;
\item locally on regularity intervals for piecewise-smooth functionals.
\end{itemize}
The only additional ingredients are regularity assumptions specific to
the functional.
\section{Interpretation}
The Universal Optimality Defect plays the role of a canonical geometric
reference quantity.
Rather than deriving independent inequalities for each entropy or
security measure, the WIS framework separates the analysis into two
independent components:
\begin{enumerate}
\item the geometry of the posterior manifold, encoded by the UOD;
\item the analytical regularity of the functional.
\end{enumerate}
This separation isolates the geometric contribution from the
functional-specific one.
\section{Scope and Limitations}
The results established in this appendix should be interpreted as
\textbf{stability estimates} rather than optimal inequalities.
In particular,
\begin{itemize}
\item the constants generally depend on the region under consideration;
\item piecewise bounds do not automatically extend across ranking transitions;
\item global constants may fail to exist near the boundary of the probability simplex.
\end{itemize}
These limitations arise from the analytical properties of the
functionals rather than from the WIS geometry itself.
\section{Outlook}
The principal contribution of Part IV is therefore not a new
inequality for a particular entropy, but a unified geometric framework
in which a broad family of information-theoretic and security
functionals can be analyzed through a common distance-like quantity.
The Universal Optimality Defect provides this common reference, while
smoothness or piecewise smoothness determines the appropriate analytical
machinery.
For a comprehensive survey of related inequalities for
\(f\)-divergences and their relationships, we refer to
\cite{sason2016fdivergence} and the references therein.
